% FAC-EN-U0036
% French source: fac_body.tex lines 1345--1351.
% Authority: PDF physical p. 36; printed p. 231 (no. 38 complete).

\medskip
\phantomsection\label{fac:II:2:no:38}
\textbf{38. Coherent algebraic sheaves.}
\phantomsection\label{fac:II:2:no:38:def:algebraic-sheaf}
Let $V$ be an algebraic variety whose sheaf of local rings is $\mathcal{O}_V$. By an \textit{algebraic sheaf} on $V$ we shall mean any sheaf of $\mathcal{O}_V$-modules, in the sense of \hyperref[fac:I:1:no:6]{no.~6}. \phantomsection\label{fac:II:2:no:38:def:algebraic-homomorphism}If $\mathcal{F}$ and $\mathcal{G}$ are two algebraic sheaves, we shall say that $\varphi:\mathcal{F}\to\mathcal{G}$ is an \textit{algebraic homomorphism}, or simply a homomorphism, if it is an $\mathcal{O}_V$-homomorphism. Recall that this is equivalent to saying that each
\[
\varphi_x:\mathcal{F}_x\longrightarrow\mathcal{G}_x
\]
is $\mathcal{O}_{x,V}$-linear and that $\varphi$ carries every local section of $\mathcal{F}$ to a local section of $\mathcal{G}$.

If $\mathcal{F}$ is an algebraic sheaf on $V$, the cohomology groups $H^q(V,\mathcal{F})$ are modules over $\Gamma(V,\mathcal{O}_V)$ (cf. \hyperref[fac:I:3:no:23]{no.~23}); in particular, they are \textit{vector spaces over $K$}.

\phantomsection\label{fac:II:2:no:38:def:coherent-algebraic-sheaf}
An algebraic sheaf $\mathcal{F}$ on $V$ is said to be \textit{coherent} if it is a coherent sheaf of $\mathcal{O}_V$-modules, in the sense of \hyperref[fac:I:2:no:12]{no.~12}. By \hyperref[fac:I:2:no:15:prop:7]{Proposition~7 of no.~15} and \hyperref[fac:II:2:no:37:prop:1]{Proposition~1 above}, such a sheaf is characterized by being locally isomorphic to the cokernel of an algebraic homomorphism
\[
\varphi:\mathcal{O}_V^q\longrightarrow\mathcal{O}_V^p.
\]

We shall give some examples of coherent algebraic sheaves; others will appear later, notably in nos.~\hyperref[fac:II:4:no:48]{48} and~\hyperref[fac:III:2:no:57]{57}.
